\subsection{Limitations}

The primary evaluation covers 14 subject groups, comprising ten from CAUCAFall and four from GMDCSA-24. This limits the range of participants and fall behaviors represented. Automatically extracted poses include tracking and localization errors, while the labeling rule excludes windows with ambiguous transition overlap. Neither threshold 0.5 nor the deployed confirmation thresholds were calibrated to a particular monitoring setting. Both unmatched-event rates describe retained video segments, so they cannot estimate notification burden during continuous care.

External Le2i video F1 lies close to the all-positive baseline of 0.8125. Class imbalance and low video specificity therefore limit the evidence for cross-dataset discrimination. Reliable participant identities are unavailable for this evaluation, leaving person-independent transfer unverified.

The hardware experiment uses a small, balanced sample of synthetic videos. It establishes local execution under network isolation, but accuracy during real continuous monitoring remains untested. Changes in illumination, camera placement, and the number of people could affect pose extraction. The ten-minute replay does not establish all-day camera reliability, and power consumption remains unmeasured. Security controls were checked functionally, without penetration testing or a formal privacy analysis.
